\section{The corrected route tree: completion and equality recovery}
\label{sec:tpc38-next-gate}

The target--coupled continuation separates into exact CRT algebra, a
congruence--completed analytic output, and an equality--recovery step.  The
last step cannot be omitted.  This section records the corrected route
tree.

\subsection{What the exact algebra removes}

The unconditional reductions prove the following.

\begin{enumerate}[label=(\roman*)]
 \item The dependence on a moving target \(N\) is exactly the ratio phase
 \(\chi(N+F)\overline{\chi(N)}\); no unidentified affine shift remains in
 the finite CRT algebra.

 \item Fixed--target character orthogonality, moving--target shifted Jacobi
 Grams, and every zero--extended character convention are explicit.

 \item For fixed \(F\), all coefficient--free target collisions are
 localized by the exact multiplicity \(d_M(F)\).  Generic determinant
 strata have no target collision.

 \item Full/proper face splitting may be reversed before estimating.
 Every nonjoint ghost then cancels identically, and the completed output is
 the joint--alias sum \(\mathscr Y_{\rm alias}\).

 \item The completion has \(O(\cK_X)\) columns \(F=kM\), with
 \(\cK_X=X^{1/400+o(1)}\) at the endpoint.  On them, every
 multiplicative character phase is one.

 \item Diagonal trace plus coefficient--blind Cauchy loses exactly the
 finite alias count \(K_B\) on the completed alias sum; at the endpoint
 this has order \(\cK_X\), and the loss is sharp in the abstract Hilbert
 column class.

 \item The original physical output is the single column
 \(\mathscr Z_0\), not \(\mathscr Y_{\rm alias}\).  The completed fixed
 fields do not recover this column in the abstract class.
\end{enumerate}

The corrected implication ledger is therefore
\begin{equation}
 \left.
 \begin{array}{c}
  \text{two fixed face--field estimates}\\
  \text{plus the exact period ledgers}
 \end{array}
 \right\}
 \Longrightarrow
 \|\mathscr Y_{\rm alias}\|_2^2
 \ll X^{o(1)}Q^2J\cK_X,
 \qquad
 \mathscr Y_{\rm alias}\not\Rightarrow\mathscr Z_0.
 \label{eq:tpc38-corrected-implication-ledger}
\end{equation}
The final nonimplication is witnessed by
\cref{thm:tpc38-equality-recovery-obstruction}.

\subsection{Route A: target--scale two--sum recovery}

The weakest explicit completion--based recovery route isolated here is to prove the completed
bound in \eqref{eq:tpc38-alias-total-target-bound} and, independently, the
nonzero--alias aggregate bound
\begin{equation}
 \left\|\sum_{\substack{k\in\cI_B\\k\ne0}}\mathscr Z_k\right\|_2^2
 \ll X^{o(1)}Q^2J K_B.
 \label{eq:tpc38-next-nonzero-alias}
\end{equation}
Their difference is \(\mathscr Z_0\), so
\cref{prop:tpc38-two-sum-recovery} gives the inherited target
\(X^{o(1)}Q^3/J\).

This route does not ask for uniform control of every alias column or for an
operator estimate under arbitrary alias weights.  It asks only for one
additional physical vector.  Its expansion contains the averaged
\(M\)-spaced M\"obius correlations in
\cref{prop:tpc38-terminal-alias-content}; hence it still requires actual
coefficient information, but it is weaker than pointwise Chowla and weaker
than a hereditary all--field theorem.

\subsection{Route B: strong localized or Gram theorems}

One may instead prove the uniform raw--column estimate
\begin{equation}
 \|\mathscr Z_k\|_2^2\ll X^{o(1)}Q^2J
 \qquad(k\in\cI_B).
 \label{eq:tpc38-next-strong-column}
\end{equation}
This is a meaningful moving--shift theorem, because the shifts
\begin{equation}
 h_k=h+kM
 \label{eq:tpc38-next-moving-h}
\end{equation}
have new primitivity and incidence behavior.  It is not merely a uniform
version of the old input--copy diagonal: at \(k=0\),
\eqref{eq:tpc38-next-strong-column} already proves the original physical
output at the stronger \(Q^2J\) scale.

The still stronger operator route is
\begin{equation}
 \|\mathsf G_B\|_{\op}\ll X^{o(1)}Q^2J.
 \label{eq:tpc38-next-strong-Gram}
\end{equation}
Since a Gram operator controls its diagonal entries,
\eqref{eq:tpc38-next-strong-Gram} implies
\eqref{eq:tpc38-next-strong-column}.  These are therefore nested
alternatives, not two independent sequential gates.  The Gram theorem is
reusable under arbitrary alias weights; the localized theorem is already
sufficient for the original output.

\subsection{Route C: exact recovery by several completions}

A different possibility is to construct several CRT completions with
moduli \(M_d\) and weights \(w_d\) such that, throughout the determinant
window,
\begin{equation}
 \sum_d w_d\one_{M_d\mid F}=\one_{F=0}.
 \label{eq:tpc38-multi-completion-recovery}
\end{equation}
If each weighted completion can be controlled without an excessive
synthesis norm, this identity recovers equality algebraically rather than
by subtracting a nonzero--alias sum.  No such well--conditioned family is
constructed here.

The four--modulus case is the simplest exact recovery:
\(M_4>X\) makes \(\one_{M_4\mid F}=\one_{F=0}\).  However, its full face
has coefficient \(1-O(J^{-1})\) at equality, so the method loops back to
the unsuppressed original coefficient.  More moduli solve recovery but do
not by themselves solve the arithmetic energy estimate.

\subsection{Three viable outcomes for the next layer}

\begin{description}[style=nextline,leftmargin=1.8em]
 \item[Positive aggregate recovery theorem.]
 Prove \eqref{eq:tpc38-next-nonzero-alias} for the actual rows and combine
 it with a completed fixed--field estimate.  This is the most economical
 route identified by the present audit.

 \item[Positive strong theorem.]
 Prove the localized column estimate or the alias--Gram estimate.  Either
 would be independently publishable; the latter is strictly stronger and
 includes the former.

 \item[Scoped obstruction or recovery no--go.]
 Within a precisely stated physical coefficient, shift, or completion
 class, prove that no target--scale recovery functional exists, or that a
 principal alias direction is unavoidable.  Such a result would rule out
 that architecture at the stated normalization, not prove a lower bound
 for the prime--pair remainder.
\end{description}

\subsection{Relation to the parity boundary}

The exact rank--one directions and terminal M\"obius correlations are
consistent with the classical parity obstruction for sieve methods
\citep{HeathBrown1982Parity,IwaniecKowalski2004}.  This paper neither
reformulates nor resolves that obstruction.  A future positive recovery
theorem would have to exploit coefficient information beyond the finite
character identities proved here.  Even then, it would first establish the
inherited energy gate; further work would still be required before any
prime--pair asymptotic.
